\section{Physics-factored scene operator and many-body neural residual}
\label{sec:physics-factored}

A trainable scene model should not ask one neural network to rediscover
self-impedance, free-space pair coupling, rigid-motion invariance, and
many-body scattering simultaneously. vNext therefore uses a cluster/factor
decomposition whose low-order terms are physical operators.

\subsection{Object and port embeddings}

Let object $i$ carry intrinsic descriptor $\gamma_i$ and pose $g_i\in\SE(3)$.
Let $E_i$ embed the local ports of object $i$ into the global port vector.
The isolated-object contribution is
\[
  Z_i^{(1)}
  =
  Z_{\rm self}(\gamma_i,\mathcal M_{\rm bg},\omega,T).
\]
In a homogeneous isotropic background it is independent of absolute pose.

\subsection{Pair interaction}

For a pair $(i,j)$ define relative pose
\[
  g_{ij}=g_i^{-1}g_j.
\]
The two-body port interaction is
\[
  Z_{ij}^{(2)}
  =
  \mathcal K_2(
    \gamma_i,\gamma_j,g_{ij},
    \mathcal M_{\rm env},\omega,T
  ).
\]
For reciprocal media, exchanging source and receiver gives the transpose
relation required by reciprocity. The pair kernel may be analytic,
tabulated/interpolated from an integral teacher, or represented by a small
equivariant neural operator.

A first scene approximation is
\[
 Z^{[2]}
 =
 \sum_i E_i^\T Z_i^{(1)}E_i
 +\sum_{i<j}
 \left(
   E_i^\T Z_{ij}^{(2)}E_j
  +E_j^\T Z_{ij}^{(2)\T}E_i
 \right).
 \label{eq:pair-expansion}
\]

\subsection{Many-body residual}

Packages, dielectric scattering, shielding, and repeated interactions create a
non-additive remainder:
\[
  \Delta Z_{\ge3}
  =
  Z_{\rm true}-Z^{[2]}.
\]
A graph neural operator predicts this residual from physics-initialized
messages. The initial edge message is not an arbitrary learned embedding; it
contains pair features such as
\[
  Z_{ij}^{(2)},\quad
  d_{ij}/L_0,\quad
  \widehat d_{ij},\quad
  \text{electrical size},\quad
  \text{material contrast}.
\]
Message passing therefore learns repeated/many-body interaction corrections
rather than first-order coupling from scratch.

\subsection{Channel-structured residual}

The residual is decomposed into reactive and dissipative parts:
\[
 \Delta Z_{\ge3}
 =
 \Delta D_{\rm abs}
 +\Delta D_{\rm out}
 +j\Delta X.
\]
The network does not output unconstrained matrices. A structure-preserving
decoder predicts factors
\[
 D_{\rm abs}
 =
 L_{\rm abs}L_{\rm abs}\adj,\qquad
 D_{\rm out}
 =
 L_{\rm out}L_{\rm out}\adj,
\]
and a reciprocal reactance $X=X^\T$. The factors may represent the total
quantity or a positive correction to a baseline, depending on whether the
chosen residual decomposition preserves positivity termwise.

If a signed correction is necessary, the network predicts a structured latent
whose \emph{final} decoded total remains PSD; it never projects an arbitrary
matrix after the fact without reporting the correction magnitude.

\subsection{Operator-level low-rank correction}

For CERTIFIED acceleration it is useful to learn a low-rank correction to the
physical integral operator or preconditioner:
\[
  A
  \approx
  A_0+U_\theta K_\theta V_\theta\adj ,
  \qquad
  r_\theta\ll N.
\]
When appropriate symmetry permits $V_\theta=U_\theta$, structured $K_\theta$
can preserve reciprocity or energy properties. If $A_0^{-1}$ is cheap, the
Woodbury identity gives
\[
 (A_0+UKV\adj)^{-1}
 =
 A_0^{-1}
 -A_0^{-1}U
 \left(K^{-1}+V\adj A_0^{-1}U\right)^{-1}
 V\adj A_0^{-1}.
\]
This creates a bridge between the FAST neural macromodel and the physical
operator used by CERTIFIED mode.

The learned operator correction is optional; a port-only residual model may be
used first. In all cases the true physical residual remains the acceptance
criterion.

\subsection{Loss-field factorization}

The same hierarchy applies to loss:
\[
  H_q(\xi)
  =
  H_q^{\rm self}(\xi)
  +H_q^{\rm pair}(\xi)
  +\Delta H_q^{\rm many}(\xi).
\]
Self loss modes are intrinsic to each conductor/material body and move with its
pose. Pair modes account for leading proximity/coupling. The learned residual
contains higher-order redistribution. The final factorized decoder enforces
local non-negativity and global power closure.

\subsection{Thermal interaction factorization}

For homogeneous or weakly structured thermal backgrounds, the transfer between
source object $i$ and observation object $j$ has a Green-kernel pair baseline.
Interface/package effects are then learned or reduced as a many-body correction.
This mirrors the electromagnetic scene decomposition and improves transfer
across variable object counts.

\subsection{Training hierarchy}

Training follows the factorization:
\begin{enumerate}
\item learn/validate isolated-object self libraries;
\item learn/validate two-object pair kernels over relative pose;
\item freeze or slowly fine-tune low-order physics and train many-body residuals;
\item jointly fine-tune decoded observables with fixed release sets;
\item calibrate residual/uncertainty heads against CERTIFIED solves.
\end{enumerate}

This hierarchy allows pair data to be reused combinatorially across many
multi-object scenes and reduces the number of expensive full-scene REFERENCE
samples.

\subsection{Failure criterion}

If the many-body residual is routinely of the same scale as the full operator
for the target scene class, the pair baseline is considered inadequate. The
response is to improve the physical cluster expansion or environment kernel,
not merely to enlarge the neural network.
